% Section 3.8, from lectures/L03/appendix/b_exercises.md.

\section{Uppgifter}
\label{c3:sec:uppgifter}

\begin{aside}
\textbf{Så arbetar du med uppgifterna.} Del~1 och del~2 räknas under lektionen: först på egen
hand, några minuter per uppgift, sedan gemensamt. Del~3 är extrauppgifter att träna vidare på
efter lektionen. De flesta går att räkna i huvudet eller med papper och penna.

\facitnotis{L03}
\end{aside}

% ------------------------------------------------------------------------------------------
\uppgiftsdel{Del 1}{Repetitionsuppgifter}

\uppgift{1.1}{Förenkling}
Förenkla följande uttryck.
\begin{delar}(3)
  \task $3(2R + 4) - 2(R - 1)$
  \task $(U + 3)^2 - 9$
  \task $\dfrac{6I^2 + 4I}{2I}$, antag $I \neq 0$
\end{delar}

\uppgift{1.2}{Faktorisering}
Faktorisera.
\begin{delar}(3)
  \task $R^2 - 25$
  \task $3U^2 + 6U$
  \task $I^2 - 8I + 16$
\end{delar}

% ------------------------------------------------------------------------------------------
\needspace{16\baselineskip}
\uppgiftsdel{Del 2}{Nytt stoff}

\uppgift{2.1}{Ohms lag: sök motståndet}
Spänningen över ett motstånd är $U = 15\,\text{V}$ och strömmen är $I = 0{,}3\,\text{A}$.

\bild[0.31\linewidth]{lectures/L03/appendix/images/2.1_circuit.png}

Ohms lag anger att $U = R \cdot I$.
\begin{delar}(2)
  \task Sätt upp ekvationen och lös för $R$.
  \task Kontrollera svaret.
\end{delar}

\uppgift{2.2}{Seriekrets: ekvation med parentes}
Två motstånd är seriekopplade. Kretsen matas med $U = 24\,\text{V}$ och strömmen är
$I = 0{,}4\,\text{A}$. Det ena motståndet är $R_2 = 40\,\Omega$, det andra är okänt.

\bild[0.41\linewidth]{lectures/L03/appendix/images/2.2_circuit.png}

Ohms lag för hela kretsen ger
\[
  U = (R_1 + R_2) \cdot I.
\]
\begin{parts}
  \item Sätt in värdena och lös ekvationen för $R_1$ \textbf{utan} att multiplicera ut parentesen.
  \item Lös ekvationen en gång till, men multiplicera ut parentesen först. Kontrollera att du får
    samma svar.
  \item Kontrollera lösningen genom att beräkna strömmen $I = \dfrac{U}{R_1 + R_2}$.
\end{parts}

\uppgift{2.3}{Spänningsdelning}
I en spänningsdelare ges utspänningen av
\[
  U_{\text{ut}} = \frac{R_2}{R_1 + R_2} \cdot U_{\text{in}}.
\]

\bild[0.44\linewidth]{lectures/L03/appendix/images/2.3_circuit.png}

Antag att $U_{\text{in}} = 9\,\text{V}$, $R_2 = 3\,\Omega$ och $U_{\text{ut}} = 3\,\text{V}$.
\begin{parts}
  \item Sätt upp ekvationen för $U_{\text{ut}}$ och lös för $R_1$.
  \item Kontrollera svaret.
\end{parts}

\uppgift{2.4}{Tidskonstant i RC-krets}
Tidskonstanten $\tau$ (tau) i en RC-krets är
\[
  \tau = R \cdot C.
\]

\bild[0.39\linewidth]{lectures/L03/appendix/images/2.4_circuit.png}

Kondensatorns kapacitans är $C = 100\,\upmu\text{F} = 100 \times 10^{-6}\,\text{F}$ och
tidskonstanten är $\tau = 0{,}05\,\text{s}$. Beräkna motståndet $R$.

\needspace{16\baselineskip}
\uppgift{2.5}{Säkert driftområde}
En transistors effektdissipation beräknas som $P = U_{\text{CE}} \cdot I_{\text{C}}$.

\bild[0.4\linewidth]{lectures/L03/appendix/images/2.5_circuit.png}

Transistorn tål maximalt $P_{\max} = 500\,\text{mW}$ och kollektorströmmen är
$I_{\text{C}} = 50\,\text{mA}$.
\begin{parts}
  \item Sätt upp en olikhet för att bestämma det maximalt tillåtna spänningsfallet
    $U_{\text{CE}}$.
  \item Lös olikheten och ange det maximalt tillåtna $U_{\text{CE}}$ i volt.
\end{parts}

\uppgift{2.6}{Belastat batteri: olikhet med teckenvändning}
Ett batteri har tomgångsspänningen $E = 12\,\text{V}$ och den inre resistansen
$R_{\text{i}} = 0{,}8\,\Omega$. När batteriet belastas med strömmen $I$ blir klämspänningen
\[
  U = E - R_{\text{i}} \cdot I.
\]

\bild[0.43\linewidth]{lectures/L03/appendix/images/2.6_circuit.png}

\begin{parts}
  \item Beräkna klämspänningen $U$ då $I = 2\,\text{A}$.
  \item Den anslutna enheten kräver minst $9\,\text{V}$ för att fungera. Sätt upp olikheten
    $U \geq 9$ och lös den för $I$.
    \begin{aside}
    \note[Tips] Tänk på vad som händer med olikhetstecknet när du dividerar med ett negativt tal.
    \end{aside}
  \item Ett annat batteri har $E = 13\,\text{V}$ och $R_{\text{i}} = 1{,}0\,\Omega$. Vid vilken ström
    ger de båda batterierna samma klämspänning? Ställ upp en ekvation med obekant i båda led och
    lös den.
  \item Vilket av batterierna ger högst klämspänning vid $I = 8\,\text{A}$?
\end{parts}

\uppgift{2.7}{Absolutbelopp i signalanalys}
En signal kan ha positiv eller negativ avvikelse från en referensspänning. Avvikelsen $\delta$
uppfyller $\abs{\delta} \leq 0{,}5\,\text{V}$.
\begin{parts}
  \item Skriv om absolut\-belopps\-olikheten som ett dubbelsidat intervall utan absolutbelopp.
  \item Om referensspänningen är $5\,\text{V}$, ange det tillåtna intervallet för den faktiska
    spänningen.
\end{parts}

\uppgift{2.8}{Strömreglering: absolutbeloppsekvation}
En strömregulator jämför den uppmätta strömmen $I$ (i ampere) med sitt börvärde och ger
felspänningen
\[
  u_{\text{e}} = 0{,}4I - 2.
\]
\begin{parts}
  \item Vid vilken ström är felspänningen noll, det vill säga vilket är börvärdet?
  \item Regulatorn larmar när felspänningen har beloppet $1{,}2\,\text{V}$. Lös ekvationen
    $\abs{0{,}4I - 2} = 1{,}2$ som två fall och ange de två strömmar där larmet triggar.
  \item Vilket strömintervall uppfyller $\abs{0{,}4I - 2} < 1{,}2$, det vill säga när larmar
    regulatorn \emph{inte}?
\end{parts}

% ------------------------------------------------------------------------------------------
\uppgiftsdel{Del 3}{Extrauppgifter}
Uppgifterna nedan är extra träning och görs med fördel på egen hand efter lektionen.

\uppgift{3.1}{Enkla linjära ekvationer}
Lös ekvationerna.
\begin{delar}(3)
  \task $x + 7 = 12$
  \task $5x = 45$
  \task $3x - 4 = 11$
  \task $\dfrac{x}{4} = 6$
  \task $8 - 2x = 0$
  \task $-3x = 18$
\end{delar}

\uppgift{3.2}{Obekant i båda led}
Lös ekvationerna och kontrollera svaret.
\begin{delar}(3)
  \task $5x - 3 = 2x + 9$
  \task $7 - 2x = x + 1$
  \task $4(x - 1) = 2(x + 3)$
  \task $3(2x + 1) = 6x + 3$
  \task $2(x + 4) = 2x + 5$
\end{delar}

\begin{aside}
\note[Tips] Två av uppgifterna beter sig annorlunda än de övriga. Fundera på vad det betyder.
\end{aside}

\uppgift{3.3}{Ekvationer med bråk}
Lös ekvationerna.
\begin{delar}(2)
  \task $\dfrac{x}{3} + \dfrac{x}{6} = 3$
  \task $\dfrac{2x - 1}{5} = 3$
  \task $\dfrac{x}{2} - \dfrac{x}{5} = 3$
  \task $\dfrac{12}{x} = 4$, antag $x \neq 0$
\end{delar}

\uppgift{3.4}{Proportionsekvationer}
\begin{parts}
  \item Lös $\dfrac{x}{8} = \dfrac{3}{4}$.
  \item Lös $\dfrac{5}{x} = \dfrac{2}{6}$.
  \item I en spänningsdelare gäller $\dfrac{U_1}{U_2} = \dfrac{R_1}{R_2}$. Beräkna $R_1$ då
    $U_1 = 4\,\text{V}$, $U_2 = 8\,\text{V}$ och $R_2 = 2\,\text{k}\Omega$.
  \item Ett motstånd har spänningen $6\,\text{V}$ vid strömmen $0{,}2\,\text{A}$. Vilken spänning
    krävs för strömmen $0{,}5\,\text{A}$? Ställ upp en proportion och lös den.
\end{parts}

\uppgift{3.5}{Lös ut en variabel ur en formel}
Lös ut den variabel som anges.
\begin{delar}(2)
  \task $U = RI$, lös ut $I$.
  \task $P = RI^2$, lös ut $R$.
  \task $R_{\text{TOT}} = R_1 + R_2$, lös ut $R_2$.
  \task $\tau = RC$, lös ut $C$.
  \task $U = E - R_{\text{i}}I$, lös ut $R_{\text{i}}$.
  \task $\eta = \dfrac{P_{\text{ut}}}{P_{\text{in}}}$, lös ut $P_{\text{in}}$.
\end{delar}

\uppgift{3.6}{Olikheter}
Lös olikheterna och beskriv lösningen som ett intervall på tallinjen.
\begin{delar}(3)
  \task $x + 3 < 8$
  \task $2x - 1 \geq 7$
  \task $-x > 4$
  \task $-2x + 6 \leq 0$
  \task $5 - 3x > 14$
\end{delar}

\uppgift{3.7}{Dubbelsidiga olikheter}
\begin{parts}
  \item Lös $2 < x + 1 < 7$.
  \item Lös $-3 \leq 2x - 1 \leq 5$.
  \item En komponent matas med $U = 5\,\text{V}$ och tål högst $P_{\max} = 0{,}25\,\text{W}$.
    Sätt upp och lös en olikhet för den tillåtna strömmen $I$.
  \item En logisk etta tolkas när insignalen ligger i intervallet
    $2{,}0\,\text{V} \leq u \leq 3{,}3\,\text{V}$. Insignalen ges av $u = 0{,}5x$ volt. Vilka
    värden på $x$ ger en logisk etta?
\end{parts}

\uppgift{3.8}{Absolutbelopp i ekvationer}
Lös ekvationerna.
\begin{delar}(3)
  \task $\abs{x} = 6$
  \task $\abs{x - 3} = 5$
  \task $\abs{2x + 1} = 7$
  \task $\abs{x + 4} = 0$
  \task $\abs{3x - 2} = -5$
\end{delar}

\uppgift{3.9}{Absolutbelopp i olikheter och toleranser}
\begin{parts}
  \item Skriv $\abs{x - 5} \leq 2$ utan absolutbelopp och lös olikheten.
  \item Lös $\abs{2x - 6} < 4$.
  \item En matningsspänning ska vara $5\,\text{V}$ med toleransen $\pm 2\,\%$. Skriv kravet som en
    absolut\-belopps\-olikhet och ange det tillåtna spänningsintervallet.
  \item En resistor är märkt $1\,\text{k}\Omega \pm 1\,\%$ och mäts till $1\,015\,\Omega$. Uppfyller
    resistorn kravet ${\abs{R - 1000} \leq 10}$?
\end{parts}

\uppgift{3.10}{Temperaturberoende resistans}
En resistors resistans varierar med temperaturen enligt
\[
  R = R_0(1 + \alpha \cdot \Delta T).
\]
Här är $R_0 = 100\,\Omega$ resistansen vid $20\,{}^{\circ}\text{C}$,
$\alpha = 0{,}004\,{}^{\circ}\text{C}^{-1}$ och $\Delta T$ temperaturändringen räknad från
$20\,{}^{\circ}\text{C}$.
\begin{parts}
  \item Beräkna $R$ vid $70\,{}^{\circ}\text{C}$.
  \item Vid vilken temperatur är $R = 110\,\Omega$?
  \item Kretsen fungerar så länge $R \leq 130\,\Omega$. Vilket temperaturintervall är tillåtet?
  \item Lös ut $\Delta T$ ur formeln allmänt.
\end{parts}
